% ============================================================
% Relatório CPU RV32I Pipeline — Arquitetura de Computadores — UFRJ
% Compilar no Overleaf com pdflatex
% Upload: este .tex + ufrj_logo_1.jpeg, ufrj_logo_2.png
% ============================================================
\documentclass[12pt, a4paper]{article}

\usepackage[utf8]{inputenc}
\usepackage[T1]{fontenc}
\usepackage[brazil]{babel}
\usepackage{lmodern}
\usepackage[top=2.5cm, bottom=2.5cm, left=3cm, right=2.5cm]{geometry}
\usepackage{setspace}
\onehalfspacing
\setlength{\parindent}{1.25cm}
\setlength{\parskip}{0pt}

\usepackage{fancyhdr}
\pagestyle{fancy}
\fancyhf{}
\fancyhead[L]{\small CPU RV32I}
\fancyhead[R]{\small Arquitetura de Computadores}
\fancyfoot[C]{\thepage}
\renewcommand{\headrulewidth}{0.4pt}

\usepackage{graphicx}
\usepackage{float}
\usepackage{caption}
\captionsetup{font=small, labelfont=bf, justification=centering}

\usepackage{booktabs}
\usepackage{tabularx}
\usepackage{array}
\usepackage[table,xcdraw]{xcolor}
\definecolor{headerblue}{RGB}{31,78,121}
\definecolor{rowalt}{RGB}{238,243,248}

\usepackage{listings}
\definecolor{codebg}{RGB}{245,245,245}
\definecolor{codecomment}{RGB}{100,100,100}
\definecolor{codekey}{RGB}{31,78,121}

\lstdefinestyle{vhdlstyle}{
  backgroundcolor=\color{codebg},
  basicstyle=\ttfamily\footnotesize,
  keywordstyle=\color{codekey}\bfseries,
  commentstyle=\color{codecomment}\itshape,
  language=VHDL, frame=single, framerule=0.4pt,
  rulecolor=\color{gray!40}, breaklines=true,
  numbers=left, numberstyle=\tiny\color{gray},
  numbersep=8pt, showstringspaces=false,
  tabsize=2, xleftmargin=1em,
  literate={ã}{{\~{a}}}1 {õ}{{\~{o}}}1 {ç}{{\c{c}}}1
           {á}{{\'a}}1  {é}{{\'e}}1   {í}{{\'i}}1
           {ó}{{\'o}}1  {ú}{{\'u}}1   {â}{{\^a}}1
           {ê}{{\^e}}1  {à}{{\`a}}1
}

\usepackage[colorlinks=true, linkcolor=headerblue, urlcolor=headerblue]{hyperref}
\usepackage{amsmath}
\usepackage{enumitem}
\usepackage{titlesec}
\usepackage{tikz}
\usetikzlibrary{positioning, arrows.meta}

\titleformat{\section}{\large\bfseries\color{headerblue}}{\thesection}{1em}{}
\titleformat{\subsection}{\normalsize\bfseries\color{headerblue!80}}{\thesubsection}{1em}{}

\newcommand{\code}[1]{\texttt{#1}}
\newcommand{\reg}[1]{\texttt{x#1}}

\begin{document}
\pagenumbering{gobble}

% ── CAPA ─────────────────────────────────────────────────────
\begin{titlepage}
  \centering
  \begin{minipage}{0.35\textwidth}
    \centering
    \includegraphics[height=3.2cm]{ufrj_logo_1.jpeg}
  \end{minipage}
  \hfill
  \begin{minipage}{0.58\textwidth}
    \centering
    \includegraphics[height=3.2cm]{ufrj_logo_2.png}
  \end{minipage}

  \vspace{3.0cm}
  {\LARGE\bfseries Arquitetura de Computadores\\[0.6em]}
  {\Huge\bfseries\color{headerblue} CPU RV32I\\[0.8em]}
  \vspace{0.4cm}
  {\large Projeto e implementa\c{c}\~{a}o de uma CPU RISC-V de 32 bits\\[0.2em]
  com pipeline de 5 est\'{a}gios em VHDL\\[0.2em]
  com forwarding, detec\c{c}\~{a}o de hazards e flush de controle}

  \vspace{3.0cm}
  \begin{tabular}{rl}
    \textbf{Alunos:}      & Luiz Cavalini, Eduardo Viana, Henrique Kezen \\
                          & e Rafael Maur\'{i}cio \\[0.5em]
    \textbf{Reposit\'{o}rio:} & \url{https://github.com/LuizCavalini/AC-CPU-RV32I} \\[0.5em]
    \textbf{Ferramenta:}  & Digital (Neemann) + GHDL \\
  \end{tabular}
  \vfill
  {\large Rio de Janeiro, \today}
\end{titlepage}

\pagenumbering{roman}
\pagestyle{plain}
\begingroup
  \setlength{\parskip}{0pt}
  \renewcommand{\baselinestretch}{0.95}\normalfont
  \tableofcontents
\endgroup
\newpage
\listoffigures
\newpage
\pagenumbering{arabic}
\pagestyle{fancy}

% ============================================================
\section{Introdu\c{c}\~{a}o}

Este relat\'{o}rio descreve o projeto e a implementa\c{c}\~{a}o de uma CPU RISC-V de
32 bits (RV32I) com pipeline de 5 est\'{a}gios em VHDL. O objetivo \'{e} construir
um processador funcional que execute o subconjunto de instru\c{c}\~{o}es inteiras
do ISA RISC-V sem modo supervisor, usando pipeline para maximizar o throughput.

A CPU implementa as seguintes instru\c{c}\~{o}es:

\begin{table}[H]
  \centering
  \caption{Instru\c{c}\~{o}es implementadas por categoria}
  \small\rowcolors{2}{rowalt}{white}
  \begin{tabular}{ll}
    \toprule
    \rowcolor{headerblue}
    \color{white}\textbf{Categoria} & \color{white}\textbf{Instru\c{c}\~{o}es} \\
    \midrule
    Aritm\'{e}tica   & \code{add, addi, auipc, sub} \\
    L\'{o}gica       & \code{and, andi, or, ori, xor, xori} \\
    Deslocamento     & \code{sll, slli, srl, srli} \\
    Mem\'{o}ria      & \code{lw, lui, sw} \\
    Salto/Desvio     & \code{jal, jalr, beq, bne} \\
    \bottomrule
  \end{tabular}
\end{table}

% ============================================================
\section{Arquitetura do Pipeline}

\subsection{Vis\~{a}o geral dos 5 est\'{a}gios}

A CPU \'{e} organizada em cinco est\'{a}gios cl\'{a}ssicos separados por registradores
de pipeline:

\begin{enumerate}[leftmargin=2em]
  \item \textbf{IF (Instruction Fetch):} O PC endere\c{c}a a mem\'{o}ria de
    instru\c{c}\~{o}es e a instru\c{c}\~{a}o \'{e} carregada no registrador IF/ID.
  \item \textbf{ID (Instruction Decode):} A instru\c{c}\~{a}o \'{e} decodificada,
    os operandos s\~{a}o lidos do banco de registradores e o imediato \'{e}
    gerado por extens\~{a}o de sinal.
  \item \textbf{EX (Execute):} A ALU executa a opera\c{c}\~{a}o. A Forwarding Unit
    seleciona os operandos corretos. Branches e jumps calculam seus alvos.
  \item \textbf{MEM (Memory Access):} \code{lw} l\^{e} e \code{sw} escreve na
    mem\'{o}ria de dados.
  \item \textbf{WB (Write Back):} O resultado final \'{e} escrito no banco de
    registradores.
\end{enumerate}

\subsection{Diagrama de blocos}

\begin{figure}[H]
  \centering
  \begin{tikzpicture}[
    stage/.style={draw, rectangle, minimum width=1.5cm, minimum height=1.0cm,
                  fill=blue!10, font=\small\bfseries},
    reg/.style={draw, rectangle, minimum width=0.45cm, minimum height=1.2cm,
                fill=gray!20, font=\tiny, inner sep=2pt},
    arr/.style={->, >=Stealth, thick},
    fwdarr/.style={->, >=Stealth, dashed, thick, color=teal!70},
    hazarr/.style={->, >=Stealth, dashed, thick, color=red!60},
  ]
    % ── Estágios e registradores (linha principal) ──────────────
    \node[stage] (IF)  at (0,0)   {IF};
    \node[reg]   (r1)  at (1.1,0) {\rotatebox{90}{\tiny IF/ID}};
    \node[stage] (ID)  at (2.2,0) {ID};
    \node[reg]   (r2)  at (3.3,0) {\rotatebox{90}{\tiny ID/EX}};
    \node[stage] (EX)  at (4.4,0) {EX};
    \node[reg]   (r3)  at (5.5,0) {\rotatebox{90}{\tiny EX/MEM}};
    \node[stage] (MEM) at (6.6,0) {MEM};
    \node[reg]   (r4)  at (7.7,0) {\rotatebox{90}{\tiny MEM/WB}};
    \node[stage] (WB)  at (8.8,0) {WB};

    % Setas da linha principal
    \foreach \a/\b in {IF/r1,r1/ID,ID/r2,r2/EX,EX/r3,r3/MEM,MEM/r4,r4/WB}
      \draw[arr] (\a) -- (\b);

    % ── Forwarding (abaixo do pipeline) ─────────────────────────
    % EX/MEM -> EX: ponto de saída = centro-baixo de r3, alvo = centro-baixo de EX
    \draw[fwdarr]
      (5.5,-0.6) -- (5.5,-1.6) -- (4.4,-1.6) -- (4.4,-0.5);
    \node[font=\tiny, color=teal!60] at (4.95,-1.85) {EX/MEM $\to$ EX};

    % MEM/WB -> EX: ponto de saída = centro-baixo de r4, alvo = EX mais abaixo
    \draw[fwdarr]
      (7.7,-0.6) -- (7.7,-2.4) -- (4.4,-2.4) -- (4.4,-0.5);
    \node[font=\tiny, color=teal!60] at (6.05,-2.65) {MEM/WB $\to$ EX};

    % ── Hazard / controle (acima do pipeline) ───────────────────
    % stall/flush: de ID/EX sobe, vai para a esquerda até IF/ID
    \draw[hazarr]
      (3.3,0.6) -- (3.3,1.4) -- (1.1,1.4) -- (1.1,0.6);
    \node[font=\tiny, color=red!60] at (2.2,1.6) {stall / flush};

    % PC redirect: de EX sobe, vai para a esquerda até IF
    \draw[hazarr]
      (4.4,0.5) -- (4.4,2.2) -- (0,2.2) -- (0,0.5);
    \node[font=\tiny, color=red!60] at (2.2,2.4) {PC redirect (branch/jump)};

  \end{tikzpicture}
  \caption{Pipeline de 5 est\'{a}gios com forwarding e detec\c{c}\~{a}o de hazards}
  \label{fig:pipeline}
\end{figure}

% ============================================================
\section{M\'{o}dulos VHDL}

\subsection{ALU (\code{src/alu/alu.vhd})}

A ALU de 32 bits recebe dois operandos e um seletor de opera\c{c}\~{a}o de 4 bits:

\begin{table}[H]
  \centering
  \caption{Opera\c{c}\~{o}es da ALU}
  \small\rowcolors{2}{rowalt}{white}
  \begin{tabular}{cll}
    \toprule
    \rowcolor{headerblue}
    \color{white}\textbf{alu\_ctrl} & \color{white}\textbf{Opera\c{c}\~{a}o} &
    \color{white}\textbf{Usada por} \\
    \midrule
    \code{0000} & ADD & \code{add, addi, lw, sw, auipc, jalr} \\
    \code{0001} & SUB & \code{sub, beq, bne} \\
    \code{0010} & AND & \code{and, andi} \\
    \code{0011} & OR  & \code{or, ori} \\
    \code{0100} & XOR & \code{xor, xori} \\
    \code{0101} & SLL & \code{sll, slli} \\
    \code{0110} & SRL & \code{srl, srli} \\
    \code{0111} & LUI & \code{lui} (passa B direto) \\
    \code{1000} & SLT & compara\c{c}\~{a}o com sinal \\
    \bottomrule
  \end{tabular}
\end{table}

A flag \code{zero\_o} \'{e} ativada quando o resultado \'{e} zero, usada por
\code{beq} (verifica \code{zero=1}) e \code{bne} (verifica \code{zero=0}).

\subsection{Banco de Registradores (\code{src/regfile/regfile.vhd})}

Armazena os 32 registradores de 32 bits. O registrador \reg{0} \'{e} fixo em zero:
escritas s\~{a}o ignoradas e leituras sempre retornam \code{0x00000000}. A escrita
\'{e} s\'{i}ncrona e a leitura \'{e} ass\'{i}ncrona (combinacional), permitindo que WB escreva
e ID leia no mesmo ciclo.

\subsection{Decodificador (\code{src/control/decoder.vhd})}

Interpreta opcode, \code{funct3} e \code{funct7} para gerar os sinais de controle.
Suporta os seis formatos RV32I (R, I, S, B, U, J) com extens\~{a}o de sinal do
imediato para cada formato.

\begin{table}[H]
  \centering
  \caption{Sinais de controle gerados pelo decodificador}
  \small\rowcolors{2}{rowalt}{white}
  \begin{tabular}{lll}
    \toprule
    \rowcolor{headerblue}
    \color{white}\textbf{Sinal} & \color{white}\textbf{Bits} &
    \color{white}\textbf{Descri\c{c}\~{a}o} \\
    \midrule
    \code{reg\_we}   & 1 & habilita escrita no banco de registradores \\
    \code{alu\_src}  & 1 & seleciona rs2 (0) ou imediato (1) \\
    \code{alu\_ctrl} & 4 & seleciona opera\c{c}\~{a}o da ALU \\
    \code{mem\_we}   & 1 & habilita escrita na mem\'{o}ria (\code{sw}) \\
    \code{mem\_re}   & 1 & habilita leitura da mem\'{o}ria (\code{lw}) \\
    \code{wb\_sel}   & 2 & fonte do writeback (ALU / MEM / PC+4) \\
    \code{branch}    & 1 & indica branch (\code{beq}/\code{bne}) \\
    \code{jump}      & 1 & indica salto (\code{jal}/\code{jalr}) \\
    \code{jalr}      & 1 & alvo = rs1+imm em vez de PC+imm \\
    \code{imm}       & 32 & imediato com extens\~{a}o de sinal \\
    \bottomrule
  \end{tabular}
\end{table}

\subsection{Registradores de Pipeline}

\begin{itemize}[leftmargin=2em]
  \item \textbf{IF/ID:} propaga PC e instru\c{c}\~{a}o. Suporta \code{flush}
    (hazard de controle) e \code{stall} (load-use).
  \item \textbf{ID/EX:} propaga todos os sinais de controle, operandos, imediato,
    endere\c{c}os de registradores e PC. Suporta \code{flush}.
  \item \textbf{EX/MEM:} propaga resultado da ALU, dado de escrita, rd e PC+4.
  \item \textbf{MEM/WB:} propaga resultado da ALU, dado lido da mem\'{o}ria, rd e
    PC+4 para writeback.
\end{itemize}

% ============================================================
\section{Tratamento de Hazards}

\subsection{Data Hazards --- Forwarding}

A \code{forwarding\_unit} resolve conflitos RAW encaminhando resultados em
tr\^{a}nsito diretamente para as entradas da ALU:

\begin{itemize}[leftmargin=2em]
  \item \textbf{EX/MEM $\to$ EX:} resultado da ALU do ciclo anterior encaminhado
    quando \code{rd\_mem = rs1\_ex} ou \code{rd\_mem = rs2\_ex}.
  \item \textbf{MEM/WB $\to$ EX:} dado de writeback encaminhado quando
    \code{rd\_wb = rs1\_ex} ou \code{rd\_wb = rs2\_ex}.
  \item \textbf{Prioridade:} EX/MEM tem prioridade sobre MEM/WB.
\end{itemize}

Codifica\c{c}\~{a}o dos muxes: \code{"00"} = sem forwarding, \code{"01"} = MEM/WB,
\code{"10"} = EX/MEM.

\subsection{Load-Use Hazard --- Stall}

Quando \code{lw} \'{e} seguido imediatamente por instru\c{c}\~{a}o que usa o valor
carregado, a \code{hazard\_unit} detecta e insere 1 ciclo de stall:

\begin{equation}
  \text{stall} = \text{mem\_re\_ex} \;\text{AND}\;
  (\text{rd\_ex} \neq \text{x0}) \;\text{AND}\;
  (\text{rd\_ex} = \text{rs1\_id} \;\text{OR}\; \text{rd\_ex} = \text{rs2\_id})
\end{equation}

A unidade mant\'{e}m PC e IF/ID inalterados e insere NOP no ID/EX.

\subsection{Control Hazards --- Flush}

Branches e jumps s\~{a}o resolvidos no est\'{a}gio EX. Quando tomados, o pipeline:
(1) redireciona o PC para o alvo, e (2) descarta as duas instru\c{c}\~{o}es incorretas
j\'{a} buscadas (flush do IF/ID e ID/EX). Custo fixo: 2 ciclos por branch tomado.

% ============================================================
\section{Simula\c{c}\~{a}o e Verifica\c{c}\~{a}o}

\subsection{Testbenches individuais}

\begin{table}[H]
  \centering
  \caption{Resultado dos testbenches por m\'{o}dulo}
  \small\rowcolors{2}{rowalt}{white}
  \begin{tabular}{lcc}
    \toprule
    \rowcolor{headerblue}
    \color{white}\textbf{M\'{o}dulo} & \color{white}\textbf{Casos de teste} &
    \color{white}\textbf{Resultado} \\
    \midrule
    \code{alu}                & 11 & Todos passaram \\
    \code{regfile}            & 5  & Todos passaram \\
    \code{decoder}            & 8  & Todos passaram \\
    \code{cpu\_single\_cycle} & programa 13 instru\c{c}\~{o}es & PC final = 0x44 \\
    \code{cpu\_pipeline}      & programa 12 instru\c{c}\~{o}es & PC final = 0x5C \\
    \bottomrule
  \end{tabular}
\end{table}

\subsection{Programa de teste do pipeline}

\begin{lstlisting}[style=vhdlstyle, language={},
  caption={Programa de teste do pipeline (assembly)}]
addi x1, x0, 10     -- x1 = 10
addi x2, x0, 3      -- x2 = 3
add  x3, x1, x2     -- x3 = 13  (forwarding EX/MEM->EX: x1 disponivel)
sub  x4, x3, x2     -- x4 = 10  (forwarding EX/MEM->EX: x3 disponivel)
addi x5, x0, 13     -- x5 = 13
beq  x3, x5, +12    -- x3==x5? sim (13==13), salta para 0x20
-- 0x18, 0x1C: SKIPPED (flush por branch tomado)
addi x7, x0, 42     -- x7 = 42  (alvo do branch)
jal  x8, +8         -- x8 = 0x28, salta para 0x2C
-- 0x28: SKIPPED (flush por JAL)
addi x10, x0, 55    -- x10 = 55 (alvo do JAL)

-- extensao: memoria, lui, jalr, logica e deslocamento
sw   x3, 0(x0)      -- RAM[0] = 13         (exercita dmem_we)
lw   x11, 0(x0)     -- x11 = RAM[0] = 13   (exercita dmem_re)
lui  x12, 0x12345   -- x12 = 0x12345000    (exercita lui)
addi x16, x0, 0x50  -- x16 = 0x50 (alvo do jalr)
jalr x13, x16, 0    -- x13 = 0x44, salta para 0x50
-- 0x44, 0x48: SKIPPED (flush por JALR)
and  x14, x1, x2    -- x14 = 10 & 3 = 2    (exercita logica)
sll  x15, x1, x2    -- x15 = 10 << 3 = 80  (exercita deslocamento)
\end{lstlisting}

\subsection{Sa\'{i}da do GHDL}

\begin{lstlisting}[language={}, basicstyle=\ttfamily\small,
  backgroundcolor=\color{codebg}, frame=single, numbers=none]
$ ghdl -r --std=08 cpu_pipeline_tb 2>&1
...(report note): CPU pipeline: programa executado com sucesso!
...(report note): PC final = 0x0000005C

$ ghdl -r --std=08 reset_tb 2>&1
...(report note): PC=0x00000014
...(report note): PC apos reset=0x00000004
\end{lstlisting}

\subsection{Valida\c{c}\~{a}o no Digital}

O circuito foi integrado no simulador Digital utilizando o componente
\textit{Arquivo externo} com backend GHDL. A estrutura do circuito consiste em:

\begin{itemize}[leftmargin=2em]
  \item \textbf{Componente \code{top}:} wrapper VHDL que instancia a
    \code{cpu\_pipeline} e exporta sinais de depura\c{c}\~{a}o adicionais.
  \item \textbf{ROM (256 palavras $\times$ 32 bits):} carregada com o programa
    de teste via arquivo \code{program.hex}.
  \item \textbf{RAM (256 palavras $\times$ 32 bits):} mem\'{o}ria de dados com
    portas separadas de leitura e escrita.
  \item \textbf{Distribuidores:} extraem os bits {[}9:2{]} dos endere\c{c}os de
    32 bits para endere\c{c}ar as mem\'{o}rias por palavra.
\end{itemize}

Os seguintes sinais de depura\c{c}\~{a}o s\~{a}o vis\'{i}veis durante a simula\c{c}\~{a}o:

\begin{table}[H]
  \centering
  \caption{Sinais de depura\c{c}\~{a}o exportados pelo componente \code{top}}
  \small\rowcolors{2}{rowalt}{white}
  \begin{tabular}{lll}
    \toprule
    \rowcolor{headerblue}
    \color{white}\textbf{Sinal} & \color{white}\textbf{Bits} &
    \color{white}\textbf{Descri\c{c}\~{a}o} \\
    \midrule
    \code{pc\_debug\_o}    & 32 & PC atual (endere\c{c}o da instru\c{c}\~{a}o em IF) \\
    \code{instr\_debug\_o} & 32 & Instru\c{c}\~{a}o sendo decodificada \\
    \code{alu\_debug\_o}   & 32 & Endere\c{c}o de mem\'{o}ria (resultado ALU no MEM) \\
    \code{stall\_debug\_o} & 1  & Indica stall ativo (PC n\~{a}o avan\c{c}a) \\
    \code{flush\_debug\_o} & 1  & Indica flush ativo (descarte de instru\c{c}\~{a}o ap\'{o}s desvio tomado) \\
    \code{imem\_addr\_o}   & 32 & Endere\c{c}o da mem\'{o}ria de instru\c{c}\~{o}es \\
    \code{dmem\_addr\_o}   & 32 & Endere\c{c}o da mem\'{o}ria de dados \\
    \code{dmem\_wdata\_o}  & 32 & Dado a escrever na mem\'{o}ria (\code{sw}) \\
    \code{dmem\_we\_o}     & 1  & Habilita\c{c}\~{a}o de escrita na mem\'{o}ria \\
    \code{dmem\_re\_o}     & 1  & Habilita\c{c}\~{a}o de leitura da mem\'{o}ria \\
    \bottomrule
  \end{tabular}
\end{table}

A simula\c{c}\~{a}o interativa no Digital confirmou os seguintes comportamentos:

\begin{table}[H]
  \centering
  \caption{Resultados da simulação interativa no Digital}
  \small\rowcolors{2}{rowalt}{white} % 1. Adicionado para alternar as cores
  \begin{tabularx}{\textwidth}{Xl}   % 2. Removido os @{} das laterais para manter o mesmo recuo da Tabela 5
    \toprule
    \rowcolor{headerblue}
    \color{white}\textbf{Teste} & \color{white}\textbf{Resultado} \\
    \midrule
    % 3. Adicionado \rowcolor{white} e removido os @{} do multicolumn
    \rowcolor{white}\multicolumn{2}{l}{\textit{Fetch e controle de fluxo}} \\
    PC avança sequencialmente (0x00 $\to$ 0x04 $\to$ 0x08 $\to$ \ldots) & Confirmado \\
    Branch tomado (\code{beq}): PC busca 0x18 e 0x1C (bolha do pipeline) e 
      redireciona para 0x20 quando o branch resolve no EX & Confirmado \\
    Flush do ID/EX no \code{beq}: as instruções buscadas em 0x18/0x1C 
      são descartadas e nunca aparecem em \code{alu\_debug} & Confirmado \\
    Salto indireto (\code{jalr}): PC busca 0x44 e 0x48 (bolha do pipeline) e 
      redireciona para 0x50 quando o jalr resolve no EX & Confirmado \\
    Flush do ID/EX no \code{jalr}: idem, para o slot de delay do salto & Confirmado \\
    Decode do \code{jalr}: \code{instr\_debug = 0x000806E7} em \code{pc\_debug = 0x40} & Confirmado \\
    Leitura da ROM: \code{instr\_debug} mostra instruções corretas & Confirmado \\
    
    \rowcolor{white}\multicolumn{2}{l}{\textit{Memória de dados}} \\
    Escrita (\code{sw}): \code{RAM[0x00] = 0x0000000D} (13) & Confirmado \\
    Leitura (\code{lw}): \code{x11 = 13} (valor lido de volta) & Inferido\\
    
    \rowcolor{white}\multicolumn{2}{l}{\textit{Reset}} \\
    Reset interativo: PC retorna a 0x00 após pulso de \code{rst\_i} & Confirmado \\
    \bottomrule
  \end{tabularx}
  \footnotetext{A leitura não foi observada diretamente pois o banco de registradores
  não é exposto como sinal de depuração; a correção é garantida pela
  combinação da escrita correta em memória (linha acima) com a validação
  unitária do banco de registradores (Tabela 4, leitura assíncrona/escrita
  síncrona).}
\end{table}

\enlargethispage{3cm}
\vspace{-0.5cm}
\begin{figure}[H]
   \centering
   \includegraphics[width=0.9\textwidth]{Captura de tela de 2026-07-05 00-01-59.png}
    \caption{Simulação da CPU RV32I no Digital — PC e sinais de depuração}
    \label{fig:digital}
\end{figure}

\textbf{Observação sobre o reset no Digital:} o reset síncrono foi
verificado tanto via testbench GHDL (\code{reset\_tb}, Seção~5.3) quanto
interativamente no simulador Digital: ao ativar \code{rst\_i} por um
ciclo durante a simulação em andamento, o PC retorna corretamente a
\code{0x00} (endereço do vetor de reset), e volta a incrementar
normalmente a partir daí (\code{0x00} $\to$ \code{0x04} $\to \ldots$)
assim que \code{rst\_i} é liberado. O valor \code{0x00000004} reportado
pelo testbench \code{reset\_tb} (Seção~5.3) corresponde a uma amostra do
PC um ciclo de clock após o término do reset --- ou seja, já incrementado
a partir de \code{0x00} pelo primeiro fetch pós-reset --- e não ao valor
imediato assumido pelo PC durante o próprio pulso de reset. Os dois
resultados são, portanto, consistentes entre si.

\subsection{Bug encontrado e corrigido: flush incompleto do ID/EX}

Durante a valida\c{c}\~{a}o interativa no Digital, observou-se que o sinal
\code{flush\_debug\_o} nunca era ativado durante um branch tomado, apesar do
PC ser redirecionado corretamente. Investiga\c{c}\~{a}o do c\'{o}digo-fonte revelou
dois problemas:

\begin{enumerate}[leftmargin=2em]
  \item O pino \code{flush\_debug\_o} no wrapper \code{top.vhd} estava
    conectado a uma heur\'{i}stica derivada de sinais de barramento externos
    (\code{dmem\_we\_o or dmem\_re\_o}), sem rela\c{c}\~{a}o com o flush real do
    pipeline.
  \item Mais grave: o sinal \code{flush\_idex\_s}, que controla o descarte
    do registrador ID/EX, era acionado apenas pela detec\c{c}\~{a}o de load-use
    hazard (\code{hazard\_unit}), nunca pelo redirecionamento de PC
    (\code{pc\_redirect\_s}) causado por um branch tomado ou jump. Isso
    significava que, embora o registrador IF/ID fosse corretamente
    limpo, a instru\c{c}\~{a}o do caminho errado que j\'{a} havia avan\c{c}ado para o
    ID/EX no mesmo ciclo do redirecionamento escapava do descarte e
    seria executada indevidamente.
\end{enumerate}

A corre\c{c}\~{a}o consistiu em: (1) exportar sinais de depura\c{c}\~{a}o reais
diretamente do \code{cpu\_pipeline} (\code{pc\_if\_s}, \code{instr\_id\_s},
\code{alu\_result\_ex\_s}, \code{stall\_pc\_s}, \code{pc\_redirect\_s}), eliminando
as heur\'{i}sticas do wrapper;\newpage e (2) alterar a l\'{o}gica de flush do ID/EX para\code{flush\_idex\_s <= flush\_idex\_haz\_s or pc\_redirect\_s}, garantindo o
descarte tanto em load-use hazard quanto em desvio tomado.

A corre\c{c}\~{a}o foi validada por um teste dirigido: instru\c{c}\~{o}es
"veneno" (\code{addi} com valores identific\'{a}veis) foram inseridas nos
slots de atraso do \code{beq} e do \code{jalr}.
Ap\'{o}s a corre\c{c}\~{a}o,
\code{alu\_debug\_o} nunca exibiu esses valores em nenhum ciclo, confirmando
que as instru\c{c}\~{o}es de caminho errado s\~{a}o corretamente descartadas.
Adicionalmente, os testbenches \code{cpu\_pipeline\_tb} e \code{reset\_tb} do
GHDL foram reexecutados ap\'{o}s a corre\c{c}\~{a}o (PC final = 0x5C e reset para
0x04, respectivamente), confirmando aus\^{e}ncia de regress\~{a}o.

% ============================================================
\section{Extensão Vetorial (Parte 2)}

\subsection{Visão Geral}

A segunda etapa do projeto estende a CPU RV32I da Parte~1 para suportar
versões vetoriais de oito instruções: \code{add, addi, auipc, sub, sll,
slli, srl, srli}. A extensão reaproveita o somador vetorial de largura de
lane configurável (\code{cla4}/\code{vec\_adder}) desenvolvido no projeto
AC-AVX-Adder, complementado por um novo deslocador vetorial
(\code{vec\_shifter}) com a mesma filosofia de contenção de lane. Todas
as 20 instruções escalares da Parte~1 permanecem suportadas sem
alteração de comportamento.

\subsection{Codificação das Instruções Vetoriais}

As oito instruções vetoriais usam o opcode \code{custom-0}
(\code{0001011}), reservado pela especificação RISC-V para extensões não
padronizadas (RISC-V Unprivileged ISA, tabela de opcodes). O campo
\code{funct3} (\code{instr[14:12]}) seleciona a operação de forma
uniforme entre os oito casos, incluindo o \code{vauipc}, que no formato
U-type padrão não possui esse campo --- aqui ele é deliberadamente
reaproveitado por uniformidade de decodificação, ao custo de parte do
imediato do \code{vauipc} (ver Tabela~\ref{tab:vecsize-format}).

\begin{table}[H]
  \centering
  \caption{Codificação \code{funct3} das instruções vetoriais}
  \small\rowcolors{2}{rowalt}{white}
  \begin{tabular}{cllc}
    \toprule
    \rowcolor{headerblue}
    \color{white}\textbf{funct3} & \color{white}\textbf{Instrução} &
    \color{white}\textbf{Formato} & \color{white}\textbf{alu\_ctrl} \\
    \midrule
    \code{000} & \code{vadd}  & R-type-like & \code{0000} \\
    \code{001} & \code{vaddi} & I-type-like & \code{0000} \\
    \code{010} & \code{vauipc}& U-type-like & \code{0000} \\
    \code{011} & \code{vsub}  & R-type-like & \code{0001} \\
    \code{100} & \code{vsll}  & R-type-like (shamt = rs2[4:0]) & \code{0101} \\
    \code{101} & \code{vslli} & I-type-like (shamt = imm[4:0]) & \code{0101} \\
    \code{110} & \code{vsrl}  & R-type-like (shamt = rs2[4:0]) & \code{0110} \\
    \code{111} & \code{vsrli} & I-type-like (shamt = imm[4:0]) & \code{0110} \\
    \bottomrule
  \end{tabular}
\end{table}

Os quatro códigos \code{alu\_ctrl} reaproveitados são exatamente os já
existentes na ALU escalar da Parte~1 (Tabela~2) --- nenhuma codificação
nova foi necessária para o controle da ALU, apenas um novo caminho de
decodificação que chega até eles.

Um campo \code{vecsize} de 2 bits (mesma convenção do \code{vec\_adder}
do AC-AVX-Adder: \code{00}=lanes de 4~bits, \code{01}=8~bits,
\code{10}=16~bits, \code{11}=32~bits) é extraído de posições distintas
conforme o formato:

\begin{table}[H]
  \centering
  \caption{Posição do campo \code{vecsize} e do imediato por formato}
  \label{tab:vecsize-format}
  \small\rowcolors{2}{rowalt}{white}
  \begin{tabular}{lll}
    \toprule
    \rowcolor{headerblue}
    \color{white}\textbf{Formato} & \color{white}\textbf{vecsize} &
    \color{white}\textbf{Imediato} \\
    \midrule
    R-type-like (\code{vadd/vsub/vsll/vsrl}) & \code{funct7[1:0]} & --- \\
    I-type-like aritmética (\code{vaddi}) & \code{imm[11:10]} &
      \code{imm[9:0]} (10 bits, reduzido de 12) \\
    I-type-like deslocamento (\code{vslli/vsrli}) & \code{imm[6:5]} &
      \code{imm[4:0]} (shamt, 5 bits, igual ao escalar) \\
    U-type-like (\code{vauipc}) & \code{instr[31:30]} &
      \code{instr[29:15]} (15 bits, reduzido de 20) \\
    \bottomrule
  \end{tabular}
\end{table}

\subsection{Módulos VHDL Novos}

\begin{itemize}[leftmargin=2em]
  \item \textbf{\code{cla4} (\code{src/vector/cla4.vhd}):} somador/
    subtrator de 4 bits com lógica de carry-lookahead, reaproveitado do
    projeto AC-AVX-Adder sem alteração de interface.
  \item \textbf{\code{vec\_adder} (\code{src/vector/vec\_adder.vhd}):}
    compõe 8 instâncias de \code{cla4} em uma cadeia de 32 bits,
    bloqueando a propagação de carry nas fronteiras de lane conforme
    \code{vecsize\_i}. Recebeu uma porta de depuração nova,
    \code{cout\_debug\_o}, exportando o carry de cada um dos 8 blocos ---
    resposta direta ao critério de avaliação do enunciado ("um somador
    pode exportar seu carry para depuração").
  \item \textbf{\code{vec\_shifter} (\code{src/vector/vec\_shifter.vhd}):}
    deslocador de barril (esquerda/direita) implementado por seleção de
    bit-fonte: para cada bit de saída, calcula-se o bit de entrada
    correspondente e força-se zero caso a fonte caia fora da lane atual
    --- mesma filosofia de contenção de lane do \code{vec\_adder}, agora
    aplicada a bits deslocados em vez de carry.
  \item \textbf{\code{vec\_alu} (\code{src/vector/vec\_alu.vhd}):}
    módulo topo que multiplexa entre \code{vec\_adder} (para
    \code{vadd/vaddi/vauipc/vsub}) e \code{vec\_shifter} (para
    \code{vsll/vslli/vsrl/vsrli}) com base no mesmo código
    \code{alu\_ctrl} da ALU escalar.
\end{itemize}

\subsection{Integração no Pipeline}

O \code{decoder} foi estendido com dois novos sinais de saída,
\code{is\_vector\_o} e \code{vecsize\_o}, através de um novo ramo no
\code{case opcode} para \code{custom-0}. O registrador \code{id\_ex\_reg}
passou a propagar ambos os sinais de ID para EX, com o mesmo
comportamento de \code{flush} (zeramento) dos demais sinais de controle
--- garantindo que uma instrução descartada por hazard de controle nunca
dispare o \code{vec\_alu} indevidamente.

No estágio EX, uma instância de \code{vec\_alu} opera em paralelo à ALU
escalar, consumindo os mesmos operandos já multiplexados
(\code{alu\_a\_muxed\_s}/\code{alu\_b\_muxed\_s}), que já contemplam
forwarding e o caso especial do \code{auipc} (PC como operando A). Um
novo sinal, \code{alu\_result\_muxed\_s}, seleciona entre o resultado
escalar e o vetorial conforme \code{is\_vector\_ex\_s}, e é esse sinal
combinado --- não mais \code{alu\_result\_ex\_s} diretamente --- que
alimenta o registrador EX/MEM e a saída de depuração
\code{alu\_debug\_o}. Como o \emph{forwarding} já opera a partir da
saída do registrador EX/MEM (não da ALU diretamente), a
\code{forwarding\_unit} e a \code{hazard\_unit} não precisaram de
nenhuma alteração --- o encaminhamento de resultados vetoriais funciona
"de graça".

Adicionalmente, o sinal \code{auipc\_id\_s} (que seleciona PC como
operando A da ALU) precisou ser estendido para também reconhecer o
\code{vauipc}, já que originalmente comparava apenas contra o opcode
escalar do \code{auipc}.

\subsection{Bugs Encontrados e Corrigidos Durante a Validação}

A construção do testbench de integração da Parte~2 revelou dois bugs
pré-existentes, ambos independentes da extensão vetorial em si.

\paragraph{Hazard RAW de distância-3 no banco de registradores.} O
relatório da Parte~1 já afirmava que "a leitura é assíncrona, permitindo
que WB escreva e ID leia no mesmo ciclo", mas a implementação não
cumpria essa promessa: a escrita em \code{regs()} só se tornava visível
à leitura combinacional um ciclo de clock depois de \code{we\_i}/
\code{wd\_i} estabilizarem. Isso deslocava a "transparência" de
distância-3 para distância-4, deixando dependências RAW de distância-3
sem cobertura --- longe demais para os caminhos de 1 e 2 ciclos da
\code{forwarding\_unit}, um ciclo cedo demais para a auto-resolução do
banco de registradores. O bug nunca havia sido detectado porque o
testbench original da Parte~1 só continha dependências de distância 1 e
2. Confirmado com um repro escalar isolado (independente de qualquer
instrução vetorial) e corrigido com um bypass clássico de
escrita-leitura no mesmo ciclo em \code{regfile.vhd}.

\paragraph{Encadeamento de carry em subtração vetorial multi-bloco.} O
\code{cla4.vhd} injetava o "+1" do complemento de dois
incondicionalmente em todo bloco sempre que \code{sub\_i='1'}
(\code{c(0) <= cin\_i or sub\_i}), em vez de fazê-lo apenas no primeiro
bloco de uma cadeia. Isso corrompia qualquer subtração vetorial com lane
maior que 4 bits sempre que o borrow real entre blocos deveria ser zero.
O bug foi isolado por rastreamento sinal-a-sinal ciclo a ciclo (todos os
sinais de controle e operandos do \code{vec\_alu} confirmados corretos
no ciclo da falha), o que restringiu a causa à lógica interna do
somador. A correção moveu a injeção do "+1" do \code{cla4} (que agora
apenas honra o \code{cin\_i} recebido, como um somador em cascata
convencional) para o \code{vec\_adder}, que injeta o "+1" apenas nas
fronteiras de início de lane, propagando o carry real sem interferência
dentro de uma lane aberta.

Ambos os bugs foram corrigidos com regressão completa (todos os
testbenches pré-existentes revalidados) e são detalhados nos commits
\code{2fcfb7a} e \code{d58572b} do repositório.

\subsection{Simulação e Verificação}

\begin{table}[H]
  \centering
  \caption{Resultado dos testbenches da extensão vetorial}
  \small\rowcolors{2}{rowalt}{white}
  \begin{tabular}{lcc}
    \toprule
    \rowcolor{headerblue}
    \color{white}\textbf{Módulo} & \color{white}\textbf{Casos de teste} &
    \color{white}\textbf{Resultado} \\
    \midrule
    \code{cla4}                & 5  & Todos passaram \\
    \code{vec\_adder}          & 8  & Todos passaram \\
    \code{vec\_shifter}        & 7  & Todos passaram \\
    \code{vec\_alu}            & 4  & Todos passaram \\
    \code{decoder} (casos vetoriais) & 8 & Todos passaram \\
    \code{cpu\_pipeline\_vector\_tb} & programa de 29 instruções & 10/10 valores corretos \\
    \bottomrule
  \end{tabular}
\end{table}

O \code{cpu\_pipeline\_vector\_tb} executa um programa misto
(construção de operandos via instruções escalares já validadas na
Parte~1, seguidas das 8 instruções vetoriais em múltiplas configurações
de \code{vecsize}) e verifica os resultados via escrita em memória
(\code{sw}), incluindo um caso de \emph{overflow} sem vazamento entre
lanes de 4 bits, um caso de dependência imediata entre duas instruções
vetoriais (\code{vadd} seguido de \code{vsub}, testando forwarding de
resultado vetorial) e um contraste entre deslocamento contido em lane e
deslocamento de largura total.

\subsection{Validação no Digital}

O componente \code{top} do Digital já expunha os cinco sinais de
depuração da Parte~1; um sexto, \code{vec\_cout\_debug\_o}, foi
adicionado para o carry vetorial. Como o mecanismo de build do Digital
(compilação sob demanda contra uma biblioteca de trabalho do GHDL, não
mais o arquivo único \code{cpu\_pipeline\_full.vhd} usado inicialmente
na Parte~1) não possuía um processo documentado e reproduzível, foi
criado \code{digital/build.sh}, testado com sucesso a partir de uma
biblioteca de trabalho vazia.

O mesmo programa de 29 instruções do testbench de GHDL foi carregado na
ROM (\code{digital/program\_vector.hex}) e executado interativamente no
Digital. Os 10 resultados na RAM conferem exatamente com os obtidos via
GHDL:

\begin{table}[H]
  \centering
  \caption{Resultados do programa vetorial confirmados na RAM (Digital)}
  \small\rowcolors{2}{rowalt}{white}
  \begin{tabular}{lll}
    \toprule
    \rowcolor{headerblue}
    \color{white}\textbf{Instrução} & \color{white}\textbf{Endereço RAM} &
    \color{white}\textbf{Valor} \\
    \midrule
    \code{vadd} (4b, overflow)      & 0x00 & 0x00000000 \\
    \code{vsub} (8b)                & 0x04 & 0x02020202 \\
    \code{vsub} (32b, forwarding)   & 0x08 & 0xFFFFFFFF \\
    \code{vslli} (4b, contido)      & 0x0C & 0x00000000 \\
    \code{vslli} (32b)              & 0x10 & 0x11111110 \\
    \code{vsll} (8b, reg-reg)       & 0x14 & 0x44444444 \\
    \code{vsrl} (8b, reg-reg)       & 0x18 & 0x04040404 \\
    \code{vauipc}                   & 0x1C & 0x0000105C \\
    \code{vsrli} (4b, contido)      & 0x20 & 0x00000000 \\
    \code{vaddi} (8b)               & 0x24 & 0x11111175 \\
    \bottomrule
  \end{tabular}
\end{table}

Os endereços \code{0x00} e \code{0x08} correspondem exatamente às
instruções que expuseram os dois bugs corrigidos na Seção~6.5, servindo
como confirmação visual das correções fora do ambiente de simulação de
linha de comando.

\subsection{Premissas e Decisões de Projeto da Extensão Vetorial}

\begin{enumerate}[leftmargin=2em]
  \item \textbf{Opcode \code{custom-0}.} Escolhido por ser explicitamente
    reservado pela especificação RISC-V para extensões não padronizadas,
    sem conflito com nenhum opcode escalar da Parte~1.

  \item \textbf{Reaproveitamento da convenção \code{vecsize} do
    AC-AVX-Adder.} Os mesmos quatro valores (4/8/16/32 bits) usados no
    somador vetorial da atividade AVX foram adotados aqui, mantendo
    consistência entre os dois projetos.

  \item \textbf{\code{funct3} uniforme mesmo para \code{vauipc}.} O
    formato U-type padrão não possui campo \code{funct3}; reaproveitar
    essa posição de bits para o \code{vauipc} mantém a lógica de
    decodificação uniforme entre os oito casos, ao custo de reduzir o
    imediato do \code{vauipc} de 20 para 15 bits.

  \item \textbf{Imediatos reduzidos.} \code{vaddi} tem imediato de 10
    bits (vs. 12 no \code{addi} escalar) e \code{vauipc} de 15 bits (vs.
    20 no \code{auipc} escalar), ambos para acomodar o campo
    \code{vecsize}. Adequado ao escopo de um programa de teste da
    disciplina.

  \item \textbf{Semântica do \code{vauipc}.} Como PC é um valor escalar
    único, "vetorizar" \code{auipc} não tem significado físico pleno
    para lanes menores que 32 bits. A implementação roteia PC e o
    imediato pelo mesmo \code{vec\_adder} usado por \code{vadd}/
    \code{vsub}, com \code{vecsize} controlando apenas onde o carry é
    quebrado --- funcionalmente consistente com a categoria aritmética,
    mesmo que o resultado com lanes $<$32 bits não tenha uso prático
    como endereço.

  \item \textbf{Reaproveitamento dos códigos \code{alu\_ctrl}.} Nenhuma
    codificação de controle de ALU nova foi criada; o \code{vec\_alu}
    usa exatamente os mesmos quatro códigos já existentes na ALU
    escalar da Parte~1, simplificando a integração no decoder.
\end{enumerate}

% ============================================================
\section{Premissas e Decis\~{o}es de Projeto}

\begin{enumerate}[leftmargin=2em]
  \item \textbf{Branch resolvido no est\'{a}gio EX.} Reutiliza a ALU (subtra\c{c}\~{a}o
    + flag zero) sem hardware adicional, ao custo de 2 ciclos por branch tomado.

  \item \textbf{Forwarding cobre todos os casos exceto load-use.} EX/MEM$\to$EX
    e MEM/WB$\to$EX cobrem depend\^{e}ncias RAW de 1 e 2 ciclos. Load-use exige
    1 stall obrigat\'{o}rio.

  \item \textbf{Mem\'{o}rias separadas (Harvard).} Conforme o enunciado, mem\'{o}rias
    de instru\c{c}\~{a}o e dados distintas, cada uma entregando uma palavra por ciclo.

  \item \textbf{Reset s\'{i}ncrono.} Todos os registradores e o PC respondem ao
    \code{rst\_i} na borda de subida, garantindo inicializa\c{c}\~{a}o determin\'{i}stica.

  \item \textbf{Estados internos exportados.} Sinais internos relevantes (flags
    zero, sinais de forwarding, stall) s\~{a}o expostos como portas para depura\c{c}\~{a}o
    no Digital, conforme exigido pelo enunciado.

  \item \textbf{AUIPC via mux no operando A.} Quando opcode = \code{0010111},
    o operando A da ALU \'{e} PC em vez de rs1, calculando PC + imm sem
    hardware adicional.

  \item \textbf{Espa\c{c}o de endere\c{c}amento limitado a 256 palavras.} ROM e RAM
    s\~{a}o endere\c{c}adas pelos bits [9:2] do barramento de 32 bits (descartando
    os 2 LSBs, por alinhamento de palavra, e os bits superiores, fixados em
    zero), resultando em 1KB de espa\c{c}o endere\c{c}\'{a}vel para cada mem\'{o}ria. Essa
    simplifica\c{c}\~{a}o \'{e} adequada ao escopo de um programa de teste da
    disciplina e n\~{a}o compromete a arquitetura da CPU, que opera sobre o
    barramento de endere\c{c}o completo de 32 bits.

  \item \textbf{CPU inativa durante carga de dados.} Como todos os elementos
    de estado da CPU (PC e registradores de pipeline) s\~{a}o atualizados
    exclusivamente na borda de subida de \code{clk\_i}, e a carga do programa
    na ROM/RAM ocorre de forma ass\'{i}ncrona pela interface do Digital sem
    gera\c{c}\~{a}o de pulsos de clock, o estado interno da CPU permanece
    inalterado durante o carregamento de dados, atendendo ao requisito do
    enunciado sem necessidade de l\'{o}gica adicional.
\end{enumerate}

% ============================================================
\section{Conclus\~{a}o}

O projeto resultou em uma CPU RV32I funcional com pipeline de 5 est\'{a}gios,
implementando as instru\c{c}\~{o}es obrigat\'{o}rias. A valida\c{c}\~{a}o em GHDL confirmou
o funcionamento correto do forwarding, detec\c{c}\~{a}o de load-use e flush de
control hazards. O PC final de 0x5C no testbench do pipeline confirma que
branches tomados e JAL foram executados corretamente com descarte das
instru\c{c}\~{o}es incorretas.

O c\'{o}digo-fonte est\'{a} dispon\'{i}vel em
\url{https://github.com/LuizCavalini/AC-CPU-RV32I}.

\newpage
% ============================================================
\section{Refer\^{e}ncias}

\begin{enumerate}[leftmargin=2em]
  \item Patterson, D. A. \& Hennessy, J. L. \textit{Computer Organization and
    Design: RISC-V Edition}, 2\textordfeminine~ed., Elsevier, 2021. Cap\'{i}tulos 4 e 5.

  \item RISC-V International. \textit{The RISC-V Instruction Set Manual,
    Volume I: Unprivileged ISA}, vers\~{a}o 20191213.
    \url{https://riscv.org/technical/specifications/}

  \item Neemann, H. \textit{Digital --- Logic Designer and Circuit Simulator}.
    \url{https://github.com/hneemann/Digital}

  \item Pedroni, V. A. \textit{Circuit Design and Simulation with VHDL},
    2\textordfeminine~ed., MIT Press, 2010.
\end{enumerate}

\end{document}